% !PRINTED 163
\weberpage{163}

%% [running head, not body text -- commented out 2026-08-07] Page furniture (running
%% head, folio, signature) is excluded from this edition by standing policy.
%% \noindent\S.\,50.\hfill Elimination aus mehreren Gleichungen.\hfill 163

%<<CONT
Anzahl der Schnittpunkte dieser drei Oberflächen an. Ist aber $T$
identisch Null, so haben die drei Flächen entweder einen Flächentheil
oder eine Curve gemeinschaftlich.

Den Grad der Endgleichung können wir auf folgendem Wege bestimmen,
der gar keine weitläufigen Betrachtungen über die Grade der
Functionen $R$, $S$, $T$ erfordert. Wir führen, um uns von
Ausnahmefällen unabhängig zu machen, homogene Variable ein, indem
wir eine vierte, $u$, hinzufügen, und also drei homogene quaternäre
Gleichungen
\weq{4}{F(x, y, z, u) = 0, \quad \Phi(x, y, z, u) = 0, \quad \Psi(x, y, z, u) = 0}
von den Ordnungen $m$, $n$, $p$ betrachten.

Nehmen wir den besonderen Fall an, dass jede dieser drei Functionen
in lauter lineare Factoren zerfällt, so ist klar, dass, wenn die
Anzahl der gemeinsamen Lösungen nicht unendlich ist. % [sic] printed period followed by lowercase "ihre"
ihre Anzahl
gleich dem Product $mnp$ sein muss; denn wir können je einen
linearen Factor von $F$, $\Phi$, $\Psi$ gleichzeitig Null setzen und
erhalten so $mnp$ Systeme von je drei linearen Gleichungen, die,
wenn sie von einander unabhängig sind, ebenso viele Werthsysteme für
die Verhältnisse $x : y : z : u$ ergeben.

Der Schluss, dass im allgemeinen Fall der Grad der Endgleichung von
(4) ebenso gross ist, ist auf den ersten Blick bedenklich. Das
Bedenken kann aber nach einer Bemerkung von \gesperrt{F. Klein}
durch folgende einfache Ueberlegung vollständig beseitigt werden.

Denken wir uns die Endgleichung aus den Gleichungen (4) zunächst
unter der Voraussetzung gebildet, dass die Co\"efficienten in den
Gleichungen (4) von einander unabhängige Variable sind, so wird die
Resultante, die sich nach Elimination von $x$, $y$ ergiebt, gewiss
nicht identisch Null, da sie schon für specielle
Co\"efficientensysteme, z.\,B. wenn $F$, $\Phi$, $\Psi$ in lineare
Factoren zerfallen, nicht identisch verschwindet, und wir erhalten
eine homogene Endgleichung für die zwei Unbekannten $z$, $u$ von
einem gewissen Grade, der zu ermitteln ist.

Specialisirt man die Co\"efficienten irgend wie, so kann zwar
möglicherweise die Endgleichung identisch befriedigt werden; wenn
aber dieser Fall nicht eintritt, so kann ihre Ordnung nicht geändert
werden (da ja alle Glieder von derselben Ordnung sind). Da nun durch
die oben angenommene Specialisirung, die sich in dem Zerfallen der
Functionen $F$, $\Phi$, $\Psi$ ausspricht, das
%CONT>>
